% page 519 du fac-simile (image p-518 du scan)
% bloc 172.7 x 237.7 mm ; marge gauche 29.3 mm
\pgc{7.620mm}{0.000mm}{
\l{\cel{52}511.}
\l{}
\l{\sou{sendro}. (\sou{zool.)}\cel{1}Genero de fishi teleostea, familio di la \textquotedbl{}perkidei\textquotedbl{},}
\l{\cel{4}kelke simila a la perki. - L.\cel{1}\sou{sander lucioperca}.-}
\l{}
\l{\sou{\textquotedbl{}sandwich\textquotedbl{}}. Loncho dina de shinko, e c., inter du pano-lonchi}
\l{\cel{4}dina, di qua la facio qua kontaktas kun la shinko esas butri-}
\l{\cel{4}zita. - E.}
\l{}
\l{\sou{sango}. (\sou{fiziol.}) Liquido qua cirkulas, per la impulsi de kordio,}
\l{\cel{4}e per arterii e veini, en la parti diversa di la korpo, por}
\l{\cel{4}durigar la vivo di ica. - EFIS.}
\l{}
\l{\sou{sanguina}. (\sou{pri temperamento}).- En qua preponderas sango. - DEFIRS.}
\l{}
\l{\sou{sanguisugo}\cel{2}I. (\sou{zool.)}\cel{2}Anelido qua sugas la sango di la animali,}
\l{\cel{4}quan onu uzas medicinale por la sango-tiro de la vaskuli kapi-}
\l{\cel{4}lara. - L.\cel{1}\sou{hirudo medicinalis}. - II. (\sou{metaf.}) (1) Persono}
\l{\cel{4}qua richigas su, detrimente de sua kunhomi, per exhaustar li.}
\l{\cel{4}(2) To quo exhaustas. - FIS.}
\l{}
\l{\sou{sanio}. (\sou{medicino}).Materio pusoza, fetida, qua ekfluas de ula ulceri.}
\l{\cel{4}- DEFIS.}
\l{}
\l{\sou{sanskrito.}\cel{2}Linguo sakra di la bramani, di qua la formo maxim ar-}
\l{\cel{4}kaika esas renkontrebla en le\cel{1}\sou{Veda}\cel{1}(2000 yari ante I.K.).}
\l{\cel{4}DEFIRS.}
\l{}
\l{\sou{santa.}\cel{2}I. (\sou{aludante Deo od irgo quo konsakresis a Deo}).Pura de}
\l{\cel{4}irgequala makulo, de sordidajo, de desnoblajo, e c. - II.}
\l{\cel{4}Dicesas pri ula kreiti, elevita til pureso supernatura : e}
\l{\cel{4}pri la personi quin la eklezio katolika agnoskas kom tala. -}
\l{\cel{4}DEFIS.}
\l{}
\l{\sou{santalo}. (\sou{bot.}) Substanco lignatra, kun odoro aromatoza. - DEFR.}
\l{}
\l{\sou{santolino}. (\sou{bot.}) Planto odoranta, ek la familio \textquotedbl{}kompozaji\textquotedbl{},}
\l{\cel{4}uzata kom kontre-vermo. - L.\cel{1}\sou{santolina}.\cel{2}FISL.}
\l{}
\l{\sou{santuario.}\cel{1}(\sou{religio}).I. La loko maxim santa en templo, en kirko,}
\l{\cel{4}interdiktita a la profani : (1) che la Judi,la parto sekreta}
\l{\cel{4}di la templo ube gardesis la archo di kontrato ; (2) che la}
\l{\cel{4}katoliki : parto di la kirko, cirkondata, maxim-multa-kaze, per}
\l{\cel{4}balustrado, ube esas la\cel{1}\sou{chef}a altaro. - II. Loko quan onu devas}
\l{\cel{4}respektar partikulare.- DEFIS.}
\l{}
\l{\sou{sapar.-}\cel{2}(\sou{trans.)}\cel{1}Exkavar sub edifico, konstrukturo, por igar lu}
\l{\cel{4}krular. - (\sou{anke metaf.}) DEFIRS.}
\l{}
\l{\sou{sapajuo}\cel{2}(zool.)\cel{2}Simio di Amerika, di qua la kaudo esas preniva.}
\l{\cel{4}- L.\cel{1}\sou{cebus.}\cel{2}DEFIR.}
\l{}
\l{sapono. (kemio). Kombinuro ek alkalio kun la acido oleinala, mar-}
\l{\cel{4}garinala, di graso, quan onu uzas por linjo-lavar, netigar, e}
\l{\cel{4}c. - EFIS.}
}
